% ============================================================
% CHAPITRE 3
% ============================================================
\chapter{Sprint 2 : Catalogue des plats}

\section{Introduction}

Ce deuxième cycle franchit un cap opérationnel majeur pour Assiette Gala Sfaxienne : la dématérialisation de la carte traiteur. Le but n'est plus seulement technique, il s'agit d'attirer l'œil du prospect avec une vitrine appétissante tout en donnant à l'équipe en cuisine un levier pour dé-publier un plat en rupture. Ce chapitre dépoussière la mécanique du module de catalogage, des allergènes jusqu'au stockage visuel des recettes.

% --------------------------------------------------
\section{Backlog du sprint}

\begin{table}[H]
    \centering
    \caption{Backlog du sprint 2}
    \label{tab:backlog_sprint2}
    \renewcommand{\arraystretch}{1.3}
    \begin{tabular}{|c|p{8cm}|c|}
        \hline
        \textbf{ID} & \textbf{User Story} & \textbf{Points} \\
        \hline
        US-07 & En tant que visiteur, je souhaite consulter le catalogue des plats. & 3 \\
        \hline
        US-08 & En tant que visiteur, je souhaite filtrer les plats par catégorie. & 3 \\
        \hline
        US-09 & En tant que visiteur, je souhaite rechercher un plat par nom. & 2 \\
        \hline
        US-10 & En tant qu'administrateur, je souhaite ajouter un plat au catalogue. & 5 \\
        \hline
        US-11 & En tant qu'administrateur, je souhaite modifier un plat existant. & 3 \\
        \hline
        US-12 & En tant qu'administrateur, je souhaite gérer les catégories. & 3 \\
        \hline
        US-13 & En tant qu'administrateur, je souhaite gérer les allergènes. & 2 \\
        \hline
    \end{tabular}
\end{table}

% --------------------------------------------------
\section{Conception}

\subsection{Diagramme de cas d'utilisation}

\begin{figure}[H]
    \centering
    \includegraphics[width=0.85\textwidth]{uc_sprint2.png}
    \caption{Diagramme de cas d'utilisation du sprint 2}
    \label{fig:uc_sprint2}
\end{figure}

\subsection{Description textuelle des cas d'utilisation}

\subsubsection{Cas d'utilisation : Consulter le catalogue}

\begin{table}[H]
    \centering
    \caption{Description textuelle --- Consulter le catalogue}
    \label{tab:desc_consulter_catalogue}
    \renewcommand{\arraystretch}{1.3}
    \begin{tabular}{|p{4cm}|p{8.5cm}|}
        \hline
        \textbf{Cas d'utilisation} & Consulter le catalogue \\
        \hline
        \textbf{Acteur principal} & Visiteur / Client \\
        \hline
        \textbf{Précondition} & Aucune \\
        \hline
        \textbf{Postcondition} & La liste des plats est affichée \\
        \hline
        \textbf{Scénario principal} &
        \begin{enumerate}
            \item L'internaute atterrit sur la grande vitrine de repas
            \item Le serveur Express délivre le lot de plats disponibles
            \item Isolement des repas par familles (entrées, suites, soupes) via le tri latéral
            \item Utilisation de la loupe pour chercher une recette précise (ex: Couscous)
            \item Observation de la fiche (Allergènes, tarif, belle photo)
        \end{enumerate} \\
        \hline
        \textbf{Scénario alternatif} &
        \begin{itemize}
            \item[4a.] Recette inexistante : avertissement poli de carence
        \end{itemize} \\
        \hline
    \end{tabular}
\end{table}

\subsubsection{Cas d'utilisation : Gérer les plats}

\begin{table}[H]
    \centering
    \caption{Description textuelle --- Gérer les plats}
    \label{tab:desc_gerer_plats}
    \renewcommand{\arraystretch}{1.3}
    \begin{tabular}{|p{4cm}|p{8.5cm}|}
        \hline
        \textbf{Cas d'utilisation} & Gérer les plats (CRUD) \\
        \hline
        \textbf{Acteur principal} & Administrateur \\
        \hline
        \textbf{Précondition} & L'administrateur est authentifié \\
        \hline
        \textbf{Postcondition} & Le catalogue est mis à jour \\
        \hline
        \textbf{Scénario principal (ajout)} &
        \begin{enumerate}
            \item M. Daoud (ou gestionnaire) ouvre la modale de fabrication
            \item Écriture de la dénomination commerciale, tarification et descriptif
            \item Coche des potentielles traces d'arachides ou de gluten
            \item Capture d'une belle assiette dressée
            \item Envoi asynchrone du cliché vers le cloud Cloudinary
            \item Enregistrement définitif du nouveau délice dans PostgreSQL
        \end{enumerate} \\
        \hline
        \textbf{Scénario alternatif} &
        \begin{itemize}
            \item[4a.] Photo trop lourde : blocage de sécurité
            \item[2a.] Oubli de tarif : le bouton d'enregistrement reste grisé
        \end{itemize} \\
        \hline
    \end{tabular}
\end{table}

\subsection{Diagrammes de séquence}

\subsubsection{Séquence : Consulter le catalogue avec filtres}

\begin{figure}[H]
    \centering
    \includegraphics[width=0.9\textwidth]{seq_catalogue.png}
    \caption{Diagramme de séquence --- Consultation du catalogue}
    \label{fig:seq_catalogue}
\end{figure}

\subsubsection{Séquence : Ajouter un plat}

\begin{figure}[H]
    \centering
    \includegraphics[width=0.9\textwidth]{seq_ajouter_plat.png}
    \caption{Diagramme de séquence --- Ajout d'un plat}
    \label{fig:seq_ajouter_plat}
\end{figure}

\subsection{Diagramme de classes}

Le diagramme de classes du service de catalogue (figure~\ref{fig:class_sprint2}) modélise les entités suivantes : \texttt{Plat}, \texttt{Category} et \texttt{Allergen}. La relation entre \texttt{Plat} et \texttt{Category} est de type plusieurs-à-un. La relation entre \texttt{Plat} et \texttt{Allergen} est de type plusieurs-à-plusieurs.

\begin{figure}[H]
    \centering
    \includegraphics[width=0.85\textwidth]{class_sprint2.png}
    \caption{Diagramme de classes du service de catalogue}
    \label{fig:class_sprint2}
\end{figure}

Le modèle \texttt{Plat} comporte les attributs :
\begin{itemize}
    \item \texttt{id} : identifiant unique (UUID)
    \item \texttt{name} : nom du plat (unique)
    \item \texttt{description} : texte descriptif du plat
    \item \texttt{price} : prix en dinars tunisiens (type décimal)
    \item \texttt{image} : URL de l'image hébergée sur Cloudinary
    \item \texttt{isAvailable} : indicateur de disponibilité
    \item \texttt{categoryId} : clé étrangère vers \texttt{Category}
    \item \texttt{allergens} : relation plusieurs-à-plusieurs vers \texttt{Allergen}
\end{itemize}

Le modèle \texttt{Category} contient :
\begin{itemize}
    \item \texttt{id} : identifiant unique
    \item \texttt{name} : nom de la catégorie (unique)
    \item \texttt{description} : description optionnelle
    \item \texttt{plats} : collection des plats de la catégorie
\end{itemize}

Le modèle \texttt{Allergen} contient :
\begin{itemize}
    \item \texttt{id} : identifiant unique
    \item \texttt{name} : nom de l'allergène (unique)
    \item \texttt{plats} : collection des plats contenant cet allergène
\end{itemize}

\subsection{Diagramme d'états-transitions}

Le diagramme d'états-transitions (figure~\ref{fig:state_sprint2}) modélise le cycle de vie d'un plat dans le catalogue, de sa création à sa suppression éventuelle.

\begin{figure}[H]
    \centering
    \includegraphics[width=0.85\textwidth]{state_sprint2_plat.png}
    \caption{Diagramme d'états-transitions --- Cycle de vie d'un plat}
    \label{fig:state_sprint2}
\end{figure}

% --------------------------------------------------
\section{Réalisation}

\subsection{Points d'entrée de l'API}

\begin{table}[H]
    \centering
    \caption{Endpoints REST du service de catalogue}
    \label{tab:endpoints_catalog}
    \renewcommand{\arraystretch}{1.2}
    \begin{tabular}{|l|p{4.5cm}|p{5.5cm}|}
        \hline
        \textbf{Méthode} & \textbf{Route} & \textbf{Description} \\
        \hline
        GET & /api/plats & Liste des plats (avec filtres et pagination) \\
        \hline
        GET & /api/plats/:id & Détails d'un plat \\
        \hline
        POST & /api/plats & Création d'un plat (administrateur) \\
        \hline
        PUT & /api/plats/:id & Modification d'un plat (administrateur) \\
        \hline
        DELETE & /api/plats/:id & Suppression d'un plat (administrateur) \\
        \hline
        GET & /api/categories & Liste des catégories \\
        \hline
        POST & /api/categories & Création d'une catégorie (administrateur) \\
        \hline
        PUT & /api/categories/:id & Modification d'une catégorie (administrateur) \\
        \hline
        DELETE & /api/categories/:id & Suppression d'une catégorie (administrateur) \\
        \hline
        GET & /api/allergens & Liste des allergènes \\
        \hline
        POST & /api/allergens & Création d'un allergène (administrateur) \\
        \hline
        PUT & /api/allergens/:id & Modification d'un allergène (administrateur) \\
        \hline
        DELETE & /api/allergens/:id & Suppression d'un allergène (administrateur) \\
        \hline
    \end{tabular}
\end{table}

% Schéma de la base de données retiré : le diagramme de classes (figure~\ref{fig:class_sprint2}) suffit à modéliser la structure des données.



\subsection{Gestion des images avec Cloudinary}

Plutôt que d'encombrer le serveur Express avec des fichiers lourds, nous avons délocalisé l'hébergement photographique chez \textbf{Cloudinary}. Cette décision s'est imposée parce qu'Assiette Gala photographie chaque assiette avant de la publier sur ses réseaux sociaux ; il fallait donc un entrepôt externe rapide et économique.

\begin{itemize}
    \item \textbf{Déportation du poids :} aucun binaire de 3 Mo ne transite dans PostgreSQL.
    \item \textbf{Recadrage automatique :} la photo brute prise au smartphone est retaillée au format carré sans intervention manuelle.
    \item \textbf{Distribution planétaire :} chaque cliché est répliqué sur un réseau de n\oe{}uds géographiques, garantissant un affichage instantané même depuis une connexion 3G.
\end{itemize}

Le convoyage d'une photo suit quatre étapes :
\begin{enumerate}
    \item Le formulaire d'administration embarque l'image en \texttt{multipart/form-data}.
    \item Le module \texttt{multer} réceptionne le flux binaire côté Node.js.
    \item Un appel asynchrone expédie le fichier vers le nuage Cloudinary.
    \item Le lien HTTPS résultant est enregistré dans la colonne \texttt{image} du plat concerné.
\end{enumerate}

\subsection{Interfaces réalisées}

\subsubsection{Catalogue public}

\begin{figure}[H]
    \centering
    \includegraphics[width=0.9\textwidth]{ui_catalogue.png}
    \caption{Interface du catalogue des plats}
    \label{fig:ui_catalogue}
\end{figure}

L'internaute découvre l'offre culinaire sous forme de tuiles illustrées : chaque rectangle affiche la photographie du mets, sa dénomination, son tarif et la famille à laquelle il appartient. Un rail latéral permet d'isoler une catégorie particulière, tandis qu'une barre de saisie libre filtre les résultats en temps réel.

\subsubsection{Détail d'un plat}

\begin{figure}[H]
    \centering
    \includegraphics[width=0.85\textwidth]{ui_plat_detail.png}
    \caption{Interface de détail d'un plat}
    \label{fig:ui_plat_detail}
\end{figure}

En cliquant sur une tuile, le visiteur accède à un écran dédié qui déplie la photo en plein cadre, déroule le descriptif rédactionnel, affiche le tarif en dinars et liste les éventuels allergènes. Un bouton pousse le plat directement dans le panier.

\subsubsection{Administration du catalogue}

\begin{figure}[H]
    \centering
    \includegraphics[width=0.9\textwidth]{ui_admin_plats.png}
    \caption{Interface d'administration du catalogue}
    \label{fig:ui_admin_plats}
\end{figure}

Côté coulisses, l'équipe de gestion dispose d'un tableau à colonnes triées par nom, prix et catégorie. L'ajout ou la retouche d'un plat s'opère via une fenêtre surgissante (modale), et le dépôt d'une image se fait par simple glissé-lâché depuis le bureau.

\subsubsection{Gestion des catégories et allergènes}

\begin{figure}[H]
    \centering
    \includegraphics[width=0.85\textwidth]{ui_admin_categories.png}
    \caption{Interface de gestion des catégories et allergènes}
    \label{fig:ui_admin_categories}
\end{figure}

% --------------------------------------------------
\section{Tests fonctionnels}

\begin{table}[H]
    \centering
    \caption{Tests fonctionnels --- Sprint 2}
    \label{tab:tests_sprint2}
    \renewcommand{\arraystretch}{1.3}
    \begin{tabular}{|c|p{3cm}|p{3.5cm}|p{2.5cm}|}
        \hline
        \textbf{N°} & \textbf{Cas de test} & \textbf{Comportement attendu} & \textbf{Résultat} \\
        \hline
        1 & Affichage du catalogue & Plats chargés avec images, noms, prix & Réussi \\
        \hline
        2 & Filtrage par catégorie & Seuls les plats de la catégorie sélectionnée apparaissent & Réussi \\
        \hline
        3 & Recherche par nom & Résultats filtrés dynamiquement & Réussi \\
        \hline
        4 & Ajout d'un plat avec image & Plat créé, image téléversée sur Cloudinary & Réussi \\
        \hline
        5 & Modification d'un plat & Informations mises à jour & Réussi \\
        \hline
        6 & Suppression d'un plat & Plat retiré du catalogue & Réussi \\
        \hline
        7 & Ajout d'une catégorie & Catégorie créée et disponible & Réussi \\
        \hline
        8 & Ajout d'un allergène et association à un plat & Allergène créé et lié au plat & Réussi \\
        \hline
        9 & Consultation du détail d'un plat & Toutes les informations affichées & Réussi \\
        \hline
        10 & Accès admin sans authentification & Redirection vers la page de connexion & Réussi \\
        \hline
    \end{tabular}
\end{table}

% --------------------------------------------------
\section{Conclusion}

Ce passage en revue du Sprint 2 signe l'ouverture officielle de la carte traiteur. En confiant le stockage des clichés à Cloudinary et la classification des repas à un module Prisma, la gestion du menu devient un jeu d'enfant. L'Assiette Gala possède dorénavant une authentique salle d'exposition culinaire. Le terrain est déblayé pour accueillir, lors du Sprint 3, la concrétisation des ventes et des devis financiers liés à ces mets.
